\clearpage\begingroup\NoSourceHead\SourceLeaf{8}
\begin{center}
{\Large [MECHANIK, MASS UND MESSEN.]}\\[1.3em]
{\large [AMTLICHE BERICHTE.]}\\[1.3em]
[I.]\\[.7em]
BERICHT ÜBER DIE DARSTELLUNG DER HANNOVERSCHEN\\
NORMALFUSSE.
\end{center}
\noindent\hspace*{\parindent}Da ein einfacher Fuss eine gar zu kleine Länge ist, um sich zu einem
Normal-Maasse zu eignen, so habe ich vorgezogen, die Länge von zwei Fuss
darzustellen, wobei zugleich der Umstand, dass diese Länge nach der gesetzlichen
Bestimmung genau 23 englischen Zollen gleich sein soll, eine nicht
unbedeutende Erleichterung gewährte.

Das englische Maass war gegeben durch einen von \textsc{Dollond} in London
1828 aus geschlagenem Messing verfertigen Yard, mit 6 überzähligen Zollen,
also zusammen \(3\frac{1}{2}\) englische Fuss haltend. Vor der Ablieferung war dieser
Maassstab durch den Capitaine \textsc{Kater} mit dem englischen Standard-Yard
verglichen, welcher bekanntlich seitdem in dem Brande des englischen Parlamentshauses
vernichtet ist, von welchem aber mehrere in grösster Schärfe
verglichene Copien noch vorhanden sind. Jene Vergleichung bezog sich aber
nur auf die ganze Länge des Yard auf dem \textsc{Dollond}schen Maassstabe, d. i. auf
das Stück zwischen dem 0 Strich und dem 36 Zoll-Strich, welches Stück
um \(\frac{27}{100{,}000}\) des englischen Zolls kleiner gefunden wurde, als der Standard-Yard.
In solchen Einhunderttausend-Theilen des englischen Zolls werde ich
auch alle in gegenwärtigem Berichte vorkommenden Angaben ausdrücken.
\HistoricalFoot{}{1\textsuperscript{*}}

\clearpage\SourceLeaf{9}\HistoricalHead{4}{L}{PHYSIK, AMTLICHE BERICHTE.}
In Beziehung auf den erwähnten Unterschied von 27 Theilen wird es
nöthig sein, hier noch eine Erläuterung beizufügen. Er beträgt etwa die Dicke
eines feinen Spinnefadens und bezeichnet zugleich ungefähr den Grad von
Genauigkeit, womit ein geschickter Künstler vermittelst des Reisserwerks ein
gegebenes Maass zu copiren vermag, während die nachherige \textls[120]{Prüfung}, oder
was dasselbe ist, die Vergleichung zweier fertig vorliegender Maasse, vermittelst
eines guten Comparators allerdings einer bedeutend grösseren Schärfe
fähig ist, ohne jedoch jemals \textls[120]{Einen einzelnen Theil} verbürgen zu können.
Bei dem \textls[120]{ersteren} Geschäfte hingegen (der mechanischen Arbeit) kann der
Künstler nur dadurch eine grössere Schärfe erreichen, dass er nicht ermüdet,
die noch nicht genügend genau befundenen Endstriche so lange wieder wegzuschleifen
und durch neue Versuche zu ersetzen, bis endlich ein Versuch
zur Zufriedenheit gelingt: allein der grosse Zeitaufwand, welchen eben die
jedesmahligen Prüfungen selbst kosten, muss natürlich hierbei eine Grenze
setzen, zumahl da in Beziehung auf die gewöhnlichen Zwecke dergleichen
kleine Unterschiede für Nichts zu achten sind, in Beziehung auf die allerfeinsten
Bedürfnisse hingegen die genaue \textls[120]{Kenntniss der Grösse} des zurückgebliebenen
Unterschiedes eben so gut ist, wie eine ganz vollkommene Gleichheit.

Wären nun die einzelnen Zölle, in welche der \textsc{Dollond}sche Yard eingetheilt
ist, unter sich vollkommen gleich, so würde zur Darstellung des genauen
Hannoverschen Doppelfusses nichts weiter erforderlich gewesen sein, als ihn
um \(17\frac{1}{4}\) Theile grösser zu machen, als die Länge von 23 solcher Zollabtheilungen,
und mir selbst würde dann bloss die Beaufsichtigung der erwähnten
Prüfungen vermittelst des Comparators obgelegen haben. Da aber jene vollkommene
Gleichheit eine ganz unzulässige Voraussetzung war, so musste ich
zuvor mich der langwierigen Arbeit unterziehen, die stattfindenden Ungleichheiten
in den einzelnen Zollen mit möglichster Schärfe auszumitteln, deren
Hauptresultate ich hierhersetze.

1. Die Entfernung zwischen dem 0 Strich und dem 4 Zoll-Strich fand
sich 25,5 Theile kleiner, die Entfernung zwischen dem 32 Zoll-Strich hingegen
und dem Endstrich, oder dem 36 Zoll-Strich, um 13,6 Theile grösser
als der neunte Theil des ganzen Yard von 0 bis 36 Zoll.

2. Der Abstand zwischen dem 0 Strich und dem 9 Zoll-Strich fand
sich 19,1 Theile kleiner, hingegen der Abstand zwischen dem 27 Zoll-Strich
\clearpage\SourceLeaf{10}\HistoricalHead{5}{R}{ÜBER DIE HANNOVERSCHEN NORMALFUSSE.}
\noindent und dem Endstrich um 11,5 Theile grösser, als der vierte Theil des ganzen
Yard.

Hieraus ergibt sich die Folge, dass

1. der Abstand zwischen dem 4 Zoll-Strich und dem 27 Zoll-Strich um
14 Theile,

2. der Abstand zwischen dem 9 Zoll-Strich und dem 32 Zoll-Strich um
5,5 Theile grösser ist, als \(\frac{23}{36}\) vom ganzen Yard, oder als 23 Zoll desselben
sein würden, wenn alle 36 vollkommen gleich wären.

Es muss mithin der wahre Hannoversche Doppelfuss um 3,25 Theile
grösser gemacht werden, als der erste Abstand (von 4\textsuperscript{z}—27\textsuperscript{z}) oder um 11,75
Theile grösser, als der zweite Abstand (von 9\textsuperscript{z}—32\textsuperscript{z}).

Diese beiden Arten, unter denen man die Wahl hatte, oder auch beide
zugleich bei den Prüfungen benutzen konnte, sind die einfachsten unter allen
den Combinationen, welche sich sonst noch hätten machen lassen.

Der Mechanikus \textsc{Meyerstein} unternahm hierauf die Anfertigung der Normal-Maassstäbe,
zuerst aus gegossenem Messing; es mussten aber dieselben
wieder cassirt werden, da sich ergab, dass sie auf ihrer Oberfläche keine so
feine Arbeit verstatteten, wie hier erforderlich war. Es wurden daher andere
aus geschlagenem Messing angefertigt, die, weil solches in der zu einem Maassstabe
nöthigen Dicke nicht vorkommt, in stärkere Stücke aus gegossenem
Messing mit einer Nuth eingefügt werden mussten; diese Einrichtung, welche
durch den Augenschein deutlicher erkannt wird, als durch eine Beschreibung,
ist übrigens ganz eben so bei dem \textsc{Dollond}schen Yard, und gewährte daher
auch den Vortheil, dass die Ausdehnung durch die Wärme bei allen drei
Maassstäben wie verhältnissmässig gleich betrachtet werden durfte, wenn man
nur gehörig dafür sorgte, dass jene bei allen Operationen sich immer in
gleicher Temperatur erhielten. Ohne diesen Umstand sorgfältigst zu berücksichtigen,
würde an harmonische Resultate bei diesen delicaten Operationen
gar nicht zu denken gewesen sein.

Der Mechanikus \textsc{Meyerstein} bemühte sich nun, obigen Bedingungen bei
der Fixirung der Endstriche so genau zu entsprechen, wie es die mechanische
Kunst vermag, und nach mehreren noch nicht ganz genügend befundenen Versuchen
gelang endlich bei jedem Maassstabe einer so genau, dass ich glaubte, dabei
stehen bleiben zu müssen. Die nachmaligen definitiven, bei jedem der beiden
\clearpage\SourceLeaf{11}\HistoricalHead{6}{L}{PHYSIK, AMTLICHE BERICHTE.}
\noindent Maassstäbe auf 40—50 Vergleichungen beruhenden Prüfungen ergaben, dass
der als zur Absendung nach Hannover bestimmte und deshalb \textls[120]{schon vorher}
mit No. I bezeichnete um 5 Theile, der andere, zur Hinterlegung bei der
Societät bestimmte und im voraus mit No. II bezeichnete um \(2\frac{1}{2}\) Theile
zu gross gerathen ist, Unterschiede, welche nach dem, was oben bemerkt
ist, als ganz verschwindend betrachtet werden dürfen. Wäre die Numerirung
nicht schon im voraus gravirt gewesen, so würde ich, da zufällig der zweite
noch etwas genauer gerathen ist, die Bestimmungen vertauscht haben. Es ist
mir übrigens kein Beispiel bekannt, wo ein Künstler bei einer ähnlichen Arbeit
eine grössere oder nur eine eben so grosse Uebereinstimmung erreicht
hätte.

Schliesslich ist noch zu bemerken, dass auf die Halbirungen und auf die
Unterabtheilung des ersten Fusses in einfache Zolle eine eben so grosse Sorgfalt
von dem Mechanikus \textsc{Meyerstein} verwandt ist,

und dass bei dem Hannoverschen Normalmaasse eben so wie bei dem
englischen die Temperatur von \(62^{0}\) \textsc{Fahrenheit} oder \(13\frac{1}{3}^{0}\) \textsc{Reaumur} als die
Normaltemperatur zu betrachten ist, d. i., dass unter dem Hannoverschen Fuss
die Hälfte derjenigen Länge verstanden wird, welche der materielle Maassstab
in jener Temperatur darstellt.

\medskip\noindent\hspace*{\parindent}Göttingen 27. Januar 1841.\hfill /: unterz. :/\quad C. F. \textsc{Gauss}.

\bigskip\noindent\hfill\begin{minipage}{.66\linewidth}
\small\setlength{\parindent}{1.2em}
\indent Die Uebereinstimmung dieser Abschrift mit dem Original
wird hiermit bescheinigt.

Hannover, den 13. April 1841.

\begin{center}gez. \textsc{Bening}, Kanzleirath.\end{center}
\end{minipage}
\vspace{2em}
\begin{center}\rule{1.25in}{.3pt}\end{center}
\clearpage\endgroup
